% Research note ONLY; not imported into main.tex.
% Formal algebra proved; uniform actual-orbit hypotheses open.
\subsection*{Induced entrance distributions and an exact return operator}
Let $I=I_{X,j}$ be a finite dyadic shell above the first-passage
threshold $a_X$. Consider a law $\mathbb P$ on initial values whose
relevant first-passage segments are finite almost surely. Write $N_I$
for the number of recorded visits, $Y_k$ for the state at the $k$-th
visit on $\{N_I\ge k\}$, and
\[
  q_k(A)=\mathbb P(N_I\ge k,\;Y_k\in A),\qquad A\subseteq I.
\]
Define the deterministic first-return map $R_I:I\to I\cup\{\dagger\}$
by following the accelerated orbit from a shell point until the first
subsequent visit to $I$ before its first passage below $a_X$; in the
absence of such a visit, put $R_I(y)=\dagger$.

\begin{proposition}[Exact deterministic return transport]
For the positive substochastic operator
\[
  (Q_I q)(A):=q\bigl(\{y\in I:R_I(y)\in A\}\bigr)
\]
one has $q_{k+1}=Q_Iq_k$ and hence $q_k=Q_I^{k-1}q_1$.
If there exist $0\le\rho<1$, $C>0$, and a finite nonnegative measure
$\nu_I$ on $I$ such that
\[
  q_1\le C\nu_I,\qquad Q_I\nu_I\le\rho\nu_I
\]
as measure inequalities, then
\[
 \mathbb P(N_I\ge k)\le C\nu_I(I)\rho^{k-1},\qquad
 \mathbb E N_I^2\le C\nu_I(I)\frac{1+\rho}{(1-\rho)^2}.
\]
\end{proposition}
\begin{proof}
The $(k+1)$-st visit is the first return after the $k$-th visit,
which is exactly the push-forward relation. The positive operator
$Q_I$ preserves inequalities; iterating gives
$q_k\le C\rho^{k-1}\nu_I$. Summation over $I$ proves the tail bound,
and summing $\mathbb E N_I^2=\sum_{k\ge1}(2k-1)\mathbb P(N_I\ge k)$
gives the claimed constant.
\end{proof}

\paragraph{Finite-time escape event.}
For $b,H\ge1$ let $G^{(b,H)}_{I,a_X}$ contain the shell points whose
actual forward orbits reach $\le a_X$ in at most $H$ accelerated steps
without recording $b$ more visits to $I$. This set is decidable for
fixed parameters. If $\mu_k$ is the conditional law of $Y_k$ given
$N_I\ge k$, then the elementary event inclusion implies
\[
  \mathbb P(N_I\ge k+b)
   \le\mathbb P(N_I\ge k)\,[1-\mu_k(G^{(b,H)}_{I,a_X})].
\]
For an auxiliary law $\nu_k$ on actual shell states and any $\eta>0$,
the two conditions
\[
 \nu_k(G^{(b,H)}_{I,a_X})\ge\eta+\varepsilon,\qquad
 d_{\mathrm{TV}}(\mu_k,\nu_k)\le\varepsilon
\]
imply a block-geometric contraction with factor $1-\eta$.
\emph{Neither the required uniform domination nor the conditional
transport estimates for the actual Collatz family have been proved.}
